\chapter{Fondements techniques}
\label{chap:fondements}

\section*{Introduction}
\addcontentsline{toc}{section}{Introduction}

Ce chapitre expose les notions sur lesquelles repose le système. Il commence par le problème que
tout le reste cherche à résoudre --- l'hallucination --- puis présente l'architecture RAG comme
réponse à ce problème, et détaille les trois briques de recherche qu'elle mobilise : embeddings,
BM25 et fusion par rang réciproque. Il se termine par la contrainte de souveraineté, qui
n'est pas une notion technique mais qui gouverne l'ensemble des choix du chapitre suivant. Les
termes introduits ici sont repris en annexe~\ref{ann:glossaire}.

\section{Grands modèles de langage et hallucination}

Un grand modèle de langage (\anglais{Large Language Model}, LLM) est un modèle statistique
entraîné à prédire la suite la plus plausible d'une séquence de texte. Cette formulation explique
à la fois sa fluidité et sa principale faiblesse : le modèle produit ce qui est
\emph{vraisemblable}, non ce qui est \emph{vrai}. Les deux coïncident souvent ; rien ne garantit
qu'ils coïncident toujours, et rien dans la sortie ne signale les cas où ils divergent.

Ce comportement porte le nom d'\textbf{hallucination}. En contexte juridique, il est
particulièrement dangereux. Sollicité sur un article de loi qu'il ne connaît pas, un modèle
produira volontiers un numéro d'article plausible accompagné d'un contenu inventé : la réponse
paraît crédible, la référence semble vérifiable, et la forme est indiscernable de celle d'une
réponse exacte. Le lecteur ne dispose d'aucun signal interne pour trier
(figure~\ref{fig:llm-vs-rag}).

\begin{figure}[htbp]
  \centering
  % Pourquoi un modèle seul ne suffit pas : la même question, deux chaînes.
%
% Les deux réponses sont alignées sur la même verticale : c'est la comparaison
% que la figure doit rendre immédiate. Routage orthogonal, aucun croisement.
\begin{tikzpicture}[
    font=\sffamily\scriptsize,
    boite/.style={draw, line width=0.7pt, rounded corners=2pt,
                  text width=24mm, minimum height=10mm, align=center,
                  inner sep=3pt},
    sortie/.style={boite, text width=31mm},
    d/.style={draw=slate,     fill=mist,        text=ink},
    t/.style={draw=inptblue,  fill=inptblue!8,  text=navy},
    m/.style={draw=leafgreen, fill=leafgreen!8, text=leafgreen!40!black},
    faux/.style={draw=anrtred,  fill=anrtred!6,   text=anrtred!70!black},
    vrai/.style={draw=leafgreen, fill=leafgreen!8, text=leafgreen!40!black},
    fleche/.style={-{Stealth[length=2.3mm,width=1.7mm]}, line width=0.7pt,
                   draw=slate, rounded corners=2.5pt},
    verdict/.style={font=\sffamily\scriptsize, align=center, text width=31mm},
    lane/.style={font=\sffamily\scriptsize\bfseries, anchor=east, align=right},
  ]

\def\XQ{0} \def\XM{3.6} \def\XG{7.2} \def\XR{10.8}
\def\YA{0} \def\YB{-3.2}

% ------------------------------------------------------- chaîne sans ancrage
\node[boite, d]  (qa)  at (\XQ,\YA) {Question du citoyen};
\node[boite, m]  (ma)  at (\XM,\YA) {Modèle de langage\\seul};
\node[sortie, faux] (ra) at (\XR,\YA) {«\,L'article 217 prévoit\\six semaines\,»};

\draw[fleche] (qa) -- (ma);
% Le vide de cette travée est le propos : rien n'est consulté entre la
% question et la réponse.
\draw[fleche] (ma) -- node[font=\sffamily\scriptsize\itshape, text=anrtred!70!black,
                           above, inner sep=3pt] {aucune source consultée} (ra);

\node[verdict, text=anrtred!70!black, anchor=north] at ($(ra.south)+(0,-2mm)$)
  {\textbf{Invérifiable.} Aucun texte ne fonde ce numéro, et rien dans la forme
   ne le signale.};

% ------------------------------------------------------- chaîne ancrée (RAG)
\node[boite, d]  (qb)  at (\XQ,\YB) {Question du citoyen};
\node[boite, t]  (rb)  at (\XM,\YB) {Recherche dans\\le corpus};
\node[boite, m]  (mb)  at (\XG,\YB) {Modèle de langage\\+ articles fournis};
\node[sortie, vrai] (rr) at (\XR,\YB) {«\,Quatorze semaines\\(article 152)\,»};

\draw[fleche] (qb) -- (rb);
\draw[fleche] (rb) -- node[font=\sffamily\scriptsize, text=slate, above,
                           inner sep=2pt] {articles} (mb);
\draw[fleche] (mb) -- (rr);

% Le corpus alimente la recherche par le bas : il n'est traversé par aucune
% autre arête.
\node[boite, d, minimum height=8mm] (corpus) at (\XM,{\YB-1.9})
  {Corpus\\\NbArticles{} articles};
\draw[fleche] (corpus) -- (rb);

\node[verdict, text=leafgreen!40!black, anchor=north] at ($(rr.south)+(0,-2mm)$)
  {\textbf{Vérifiable.} Le numéro cité renvoie à un article réellement fourni
   au modèle.};

% ------------------------------------------------------------- libellés de voie
\node[lane, text=anrtred!70!black]   at ($(qa.west)+(-4mm,0)$) {SANS\\RAG};
\node[lane, text=leafgreen!40!black] at ($(qb.west)+(-4mm,0)$) {AVEC\\RAG};

% Séparation des deux voies
\draw[rule, line width=0.5pt, dash pattern=on 2.5pt off 2pt]
  ($(qa.south west)+(-14mm,-11mm)$) -- ($(ra.south east)+(0,-11mm)$);

\end{tikzpicture}

  \caption[Modèle seul contre modèle ancré]{La même question, deux chaînes. Le modèle seul
  produit une référence invérifiable ; le modèle ancré cite un numéro rattaché à un texte
  réellement fourni, que le lecteur peut aller relire. L'ancrage ne rend pas le modèle plus
  savant : il rend sa réponse \emph{contrôlable}.}
  \label{fig:llm-vs-rag}
\end{figure}

Deux stratégies existent pour ancrer un modèle dans des faits réels.

\begin{itemize}
  \item Le \textbf{\anglais{fine-tuning}} réentraîne le modèle sur un corpus spécialisé. Il est
        coûteux, difficile à mettre à jour à chaque modification du texte de loi, et surtout : il
        ne permet toujours pas de citer une source vérifiable, puisque la connaissance reste
        diluée dans les poids.
  \item La \textbf{génération augmentée par la recherche} (RAG) fournit au modèle les documents
        pertinents \emph{au moment de la question}, et lui demande de s'appuyer uniquement sur
        eux.
\end{itemize}

C'est la seconde qui a été retenue. La raison n'est pas d'abord une question de coût : c'est que
le RAG rend la \textbf{citation} possible, donc la \textbf{vérification} réalisable, ce qui est la
condition même de la fiabilité recherchée. Un système dont on ne peut pas contrôler les
affirmations ne peut pas être qualifié de fiable, quelle que soit la qualité apparente de ses
réponses.

\section{L'architecture RAG}

Le principe du RAG consiste à séparer deux préoccupations habituellement confondues :
\emph{savoir} et \emph{formuler}. Le savoir est stocké dans un corpus documentaire externe et
interrogeable ; le modèle de langage n'est plus qu'un moteur de reformulation appliqué à des
documents qu'on lui fournit.

Cette séparation se traduit par deux pipelines : un pipeline d'\textbf{indexation}, exécuté une
seule fois hors ligne, et un pipeline de \textbf{requête}, exécuté à chaque question. Elle a une
conséquence pratique décisive pour le débogage : il devient possible d'inspecter les documents
retrouvés \emph{avant} l'appel au modèle. Or la majorité des mauvaises réponses en RAG proviennent
d'une mauvaise \emph{recherche}, non d'une mauvaise \emph{génération} --- et sans cette
séparation, les deux causes sont indiscernables depuis la sortie.

\section{Recherche sémantique par embeddings}

Un \textbf{embedding} est la représentation d'un texte sous forme de vecteur numérique de
dimension fixe, construite de telle sorte que deux textes de sens proche produisent des vecteurs
proches dans l'espace. La proximité est mesurée par la \textbf{similarité cosinus} :

\begin{equation}
  \mathrm{sim}(\vec{a}, \vec{b}) = \frac{\vec{a} \cdot \vec{b}}{\|\vec{a}\|\,\|\vec{b}\|}
  \label{eq:cosinus}
\end{equation}

Pour des vecteurs normalisés, cette expression se réduit à un simple produit scalaire, borné
entre $-1$ et $1$. C'est la raison pour laquelle les vecteurs sont normalisés dès l'indexation
(chapitre~\ref{chap:realisation}) : la mesure devient indépendante de la métrique de distance
configurée dans la base.

L'intérêt de cette approche est qu'elle capture le \emph{sens} et non les mots. Une question posée
en langage courant --- «~mon patron peut me virer sans raison~?~» --- peut retrouver un article
rédigé en langage juridique formel sans partager avec lui le moindre terme. Sa limite est
exactement symétrique : elle est moins précise lorsque le terme exact compte, par exemple un
numéro d'article ou une expression figée comme «~Loi~65-99~».

\section{Recherche lexicale : BM25}

BM25 (\anglais{Best Matching 25}) est une fonction de classement lexicale qui score un document
selon les termes qu'il partage avec la requête. Elle pondère chaque terme par sa rareté dans le
corpus --- un mot rare est plus discriminant qu'un mot fréquent --- et normalise par la longueur
du document, afin qu'un texte long ne soit pas favorisé par sa seule taille.

BM25 est le complément exact des embeddings : excellent sur les correspondances littérales,
aveugle à la reformulation. Un système qui n'utiliserait que l'un des deux hérite mécaniquement de
sa faiblesse.

\section{Fusion des classements : \anglais{Reciprocal Rank Fusion}}

Combiner les deux méthodes pose un problème pratique immédiat : leurs scores ne sont pas
comparables. Un score BM25 n'est pas borné ; une similarité cosinus l'est entre $-1$ et $1$. Faire
la moyenne de $8$ et de $0{,}7$ n'a aucun sens.

La \textbf{fusion par rang réciproque} (\anglais{Reciprocal Rank Fusion}, RRF) contourne le
problème en ignorant les scores et en n'utilisant que le \emph{rang} de chaque document dans
chaque liste :

\begin{equation}
  \mathrm{RRF}(d) = \sum_{m \in M} \frac{1}{k + \mathrm{rang}_m(d)}
  \label{eq:rrf}
\end{equation}

où $M$ est l'ensemble des méthodes de recherche, $\mathrm{rang}_m(d)$ le rang du document $d$
selon la méthode $m$, et $k$ une constante d'amortissement dont la valeur usuelle est~60. Un
document bien classé par les deux méthodes obtient un score élevé, sans qu'aucune normalisation
d'échelle ne soit nécessaire.

Cette élégance a une contrepartie qu'il faut énoncer maintenant : en ne retenant que le rang, RRF
traite les deux méthodes comme également crédibles. Un document que \emph{seule} la méthode la
mieux informée a correctement classé ne reçoit qu'une contribution, là où deux documents
médiocrement classés par les \emph{deux} méthodes en cumulent chacun deux. La fusion peut donc
\emph{dégrader} un bon classement (figure~\ref{fig:rrf}).

\begin{figure}[htbp]
  \centering
  % Mécanique de la fusion RRF, et l'effet de dilution qu'elle produit quand
% une seule des deux méthodes trouve le bon article.
% Exemple construit pour l'exposé ; les rangs sont ceux qu'illustre le cas
% « durée minimale du repos hebdomadaire » (article attendu : 205).
\begin{tikzpicture}[font=\sffamily\scriptsize,
    ligne/.style={draw=rule, line width=0.5pt, rounded corners=1.5pt,
                  minimum height=6mm, text width=24mm, align=left,
                  inner xsep=4pt, font=\sffamily\scriptsize, fill=white},
    cible/.style={ligne, draw=leafgreen, fill=leafgreen!8, text=leafgreen!40!black,
                  font=\sffamily\scriptsize\bfseries},
    perdu/.style={ligne, draw=anrtred, fill=anrtred!5, text=anrtred!70!black,
                  font=\sffamily\scriptsize\bfseries},
    entete/.style={font=\sffamily\scriptsize\bfseries, text=navy, align=center},
    rang/.style={font=\sffamily\scriptsize, text=slate, anchor=east}]

\def\ya{0} \def\yb{-0.75} \def\yc{-1.5}

% ------------------------------------------------------- colonne 1 : dense
\node[entete] at (0,0.9) {Classement dense};
\node[cible] (d1) at (0,\ya) {art. 205};
\node[ligne] (d2) at (0,\yb) {art. 213};
\node[ligne] (d3) at (0,\yc) {art. 206};
\foreach \i/\y in {1/\ya, 2/\yb, 3/\yc}
  \node[rang] at (-1.42,\y) {\i.};

% -------------------------------------------------------- colonne 2 : BM25
\node[entete] at (4.6,0.9) {Classement BM25};
\node[ligne] (b1) at (4.6,\ya) {art. 206};
\node[ligne] (b2) at (4.6,\yb) {art. 213};
\node[perdu] (b3) at (4.6,\yc) {art. 205 \;absent};
\foreach \i/\y in {1/\ya, 2/\yb}
  \node[rang] at (3.18,\y) {\i.};
\node[rang, text=anrtred] at (3.18,\yc) {--};

% ------------------------------------------------------ colonne 3 : fusion
\node[entete] at (9.6,0.9) {Après fusion RRF};
\node[ligne] (f1) at (9.6,\ya) {art. 206 \hfill \tiny 0,0328};
\node[ligne] (f2) at (9.6,\yb) {art. 213 \hfill \tiny 0,0328};
\node[perdu] (f3) at (9.6,\yc) {art. 205 \hfill \tiny 0,0167};
\foreach \i/\y in {1/\ya, 2/\yb, 3/\yc}
  \node[rang] at (8.18,\y) {\i.};

% ------------------------------------------------------------------ liaisons
\draw[flux] (d1.east) to[out=0,in=180] (b3.west);
\draw[flux] (d3.east) to[out=0,in=180] (b1.west);
\draw[flux] (b1.east) to[out=0,in=180] (f1.west);
\draw[flux] (b2.east) to[out=0,in=180] (f2.west);
\draw[rejet] (b3.east) to[out=0,in=180] (f3.west);

% ------------------------------------------------------------------- formule
\node[anchor=north, font=\small, text=navy] at (4.6,-2.5)
  {$\displaystyle \mathrm{RRF}(d)=\sum_{m}\frac{1}{k+\mathrm{rang}_m(d)}
    \qquad k=60$};

\node[annot, anchor=north, text width=0.86\linewidth, align=left] at (4.8,-3.6)
  {L'article~205 est classé \textbf{premier} par la recherche dense, mais
   n'apparaît pas dans les candidats BM25 : il ne reçoit le crédit que d'une
   seule méthode. Les articles 206 et 213, moyennement classés par les
   \emph{deux}, cumulent deux contributions et le dépassent. La fusion a
   \textbf{dégradé} un classement dense pourtant correct.};

\end{tikzpicture}

  \caption[Mécanique de la fusion RRF]{Mécanique de la fusion RRF et son effet de dilution. La
  formule ne connaît que les rangs : elle ne peut pas savoir laquelle des deux méthodes est la
  mieux informée sur une question donnée.}
  \label{fig:rrf}
\end{figure}

Cet effet est réel, mais il est \emph{conditionnel} : il ne se manifeste que lorsqu'une des deux
méthodes domine largement l'autre sur toute la population de questions. Le
chapitre~\ref{chap:evaluation} montre que cette condition dépend entièrement du type de questions
posées --- et qu'elle n'est pas vérifiée dès que le jeu de test cesse d'être homogène.

\section{La contrainte d'exécution locale}

L'exigence de souveraineté a une conséquence technique directe : \textbf{aucun composant du cœur
du système ne peut être un service en ligne}. Cela exclut les embeddings accessibles par
interface de programmation, les services de reclassement commerciaux et les modèles de génération
propriétaires disponibles uniquement à distance. Toute la chaîne repose donc sur des modèles
libres, exécutables sur du matériel ordinaire.

Cette contrainte, loin d'être une simple limitation, structure l'ensemble des choix du chapitre
suivant : beaucoup de questions de la forme «~pourquoi pas~X~?~» s'y résument à «~X est un service
distant~; la question du citoyen ne doit pas quitter la machine~».

\section*{Conclusion}
\addcontentsline{toc}{section}{Conclusion}

L'hallucination n'est pas un défaut accidentel des modèles de langage : c'est la conséquence
directe de ce qu'ils optimisent. Le RAG n'élimine pas ce comportement, mais il en change le
statut : en fournissant les textes au moment de la question, il rend la citation possible et donc
la vérification réalisable. Les trois briques de recherche présentées ici --- embeddings, BM25,
fusion RRF --- sont complémentaires par construction, et l'une d'elles porte déjà en germe un
résultat contre-intuitif que le chapitre~\ref{chap:evaluation} établira par la mesure.
